\chapter[Formalisation des évènements]{Formalisation des évènements}
\label{cha:formalisation-evenements}

À ce jour, les noyaux des modeleurs 3D disponibles proposent un suivi de
l'évolution des entités topologiques dans un objet. %
Ces évolutions sont représentées par les changements topologiques, ou plus
simplement évènements, que les opérations de modélisation appliques aux
différentes entités. %
Ces évènements, qui diffèrent selon les outils, peuvent être catégorisés avec
une certaine précision. %
Par exemple, pour OpenCascade, les évènements suivis sont la génération, la
modification et la destruction d'une entité. %
\todo{ajouter l'exemple d'un autre noyau où les évènements seraient plus précis,
  i.e. création, scission, fusion…}

Pour effectuer le suivi des entités et enregistrer les évènements, l'ensemble de
ces outils ont recours~: (1) soit un parcours total de l'objet pour mettre
chaque entité à jour~; (2) soit à une annotation explicite des opérations de
modélisation avec des évènements sur des types d'entités topologiques
précises. %
Dans le premier cas, l'avantage est qu'il est exhaustif, mais le coût en temps
d'exécution de cette procédure dépend de la complexité de l'objet modélisé. %
Dans le second cas, le coût en temps d'exécution est instantané, mais des
erreurs peuvent être introduites lors de la conception de l'opération,
c'est-à-dire référencer le mauvais évènement pour la mauvaise topologie, voire
oublier des évènements. %

Dans ce chapitre, nous présentons notre approche qui consiste à utiliser les
règles Jerboa pour représenter les opérations de modélisation géométrique. %
Ainsi, nous tirons avantage de ce formalisme dont l'analyse syntaxique nous
permet de détecter les évènements appliqués à chaque entité topologique
représentée dans une règle. %
Dans la première section, nous présentons la liste exhaustive des évènements que
nous proposons de suivre et les formalisons dans le contexte des G-cartes. %
Dans la deuxième section, nous formalisons ces évènements dans le contexte des
règles de transformation de graphe de sorte à caractériser les conditions
suivant lesquelles chaque évènement est détecté. %
Dans la troisième section, nous étendons la détection des évènements aux règles
Jerboa notamment avec le concept d'origine d'une entité. %
Enfin, dans une quatrième section, nous présentons des exemples de détection
d'évènements. %



%%% Local Variables:
%%% mode: LaTeX
%%% TeX-master: "../../../main"
%%% End:
